%% P13 --- Grafy, DAG-i i bladzenie losowe
%%
%% Kazda liczba w tym programie, ktorej czytelnik nie policzy w pamieci,
%% pochodzi z code/p13_graphs.py i wchodzi na strone przez \val{}. Listing
%% konsolowy to zatwierdzony transcript pisany przez ten skrypt.
%%
%% TEZA: graf to zbior plus relacja; macierz sasiedztwa zamienia kazde pytanie
%% o sciezki w pytanie o potegi macierzy; a porzadek topologiczny DAG-a jest
%% tym, co czyni obliczenie dobrze okreslonym -- wiec przeplyw agenta, system
%% budowania i siec neuronowa sa jednym obiektem.
%%
%% CO TEN PROGRAM DOSTAJE, odczytane z plikow, a nie z pamieci:
%%   F10  ODDAJE "graf jako zbior plus relacja" tutaj Z NAZWY, we wlasnym
%%        naglowku pliku, i daje zbiory, regule sumy oraz liczbe par.
%%   P12  ma argumenty zliczajace i celowo nie uzyl ich na grafach. Jego
%%        przyklad z dokumentami i parami jest tu kontynuowany.
%%   P06  daje twierdzenie, na ktorym stoi sekcja 3: iloczyn macierzy JEST
%%        zlozeniem, a wiersz-razy-kolumna z tego wypada. Wlasnie dlatego
%%        A^k liczy spacery dlugosci k.
%%   P10  daje wektory wlasne. Rozklad stacjonarny to wektor wlasny dla
%%        wartosci 1, czyli ten sam rachunek pod inna nazwa.
%%   P03  daje slownik kosztu potrzebny sekcji 2.
%%
%% CO TEN PROGRAM ODDAJE DALEJ: P16 deklaruje P13 jako zaleznosc i potrzebuje
%% grafu obliczen zdefiniowanego porzadnie, jeden program przed autodiff.
%%
%% CZEGO NIE RUSZA, sprawdzone w tools/programs.json:
%%   jakobian, tryb prosty i odwrotny, co liczy backward()          -> P16
%%   kwantyfikatory, indukcja, pisanie dowodu                       -> P14
%%   czym JEST prawdopodobienstwo i co robi warunkowanie            -> P23
%%
%% Blizniak: programs/en/P13-graphs-dags.tex. Ramka w ramke, makro w makro, liczba
%% w liczbe; porownuje je tools/parity.py.

\program{Grafy, DAG-i i błądzenie losowe}
\label{prog:P13}
\index{graf}
\index{graf!macierz sąsiedztwa}

\begin{outcomes}
\outcome{1--8}{powiedzieć, czym jest graf, bez rysunku, i policzyć jego
  krawędzie wobec par, z których są brane}
\outcome{9--13}{wybrać między macierzą sąsiedztwa a listą sąsiedztwa na
  podstawie gęstości i operacji, z rachunkiem}
\outcome{14--20}{czytać potęgę macierzy sąsiedztwa jako liczbę spacerów i
  powiedzieć, czym to czyni głębokość sieci przekazującej komunikaty}
\outcome{21--26}{powiedzieć, czym jest porządek topologiczny, dlaczego zwykle
  jest ich więcej niż jeden i dlaczego wszystkie dają tę samą odpowiedź}
\outcome{27--34}{czytać rozkład stacjonarny jako wektor własny i powiedzieć,
  co robi z nim wierzchołek bez krawędzi wychodzącej}
\end{outcomes}

\begin{quiz}
\item Graf ma $\val{p13.g.ne}$ krawędzi. Ile wynosi suma jego stopni?
  \teachesat{4--5}
  \answerto{$\val{p13.g.degsum}$, dwa razy tyle co krawędzi, bo każda krawędź
  jest liczona raz przy każdym ze swoich dwóch końców. To podwójne zliczanie
  z~Programu~\ref{prog:P12} w nowym przebraniu.}
\item $\val{p13.big.n}$ wierzchołków i $\val{p13.big.m}$ krawędzi. Ile mniej
  więcej komórek ma macierz sąsiedztwa? \teachesat{9--10}
  \answerto{$\val{p13.big.cells}$ \dash{} milion do kwadratu \dash{} wobec
  $\val{p13.big.slots}$ miejsc dla listy, czyli współczynnik
  $\val{p13.big.ratio}$. Prawdziwe grafy są rzadkie, a macierz przechowuje
  się tylko wtedy, gdy zamierza się przez nią mnożyć.}
\item Co liczy wyraz $(i,j)$ macierzy $A^{2}$? \teachesat{14--15}
  \answerto{Spacery długości dwa z $i$ do $j$. Suma
  $\sum_k A_{ik}A_{kj}$ ma wyraz równy $1$ dla każdego wierzchołka pośredniego
  połączonego z obydwoma, czyli złożenie z~Programu~\ref{prog:P06} robiące
  arytmetykę.}
\item Ile porządków topologicznych ma graf skierowany bez cyklu?
  \teachesat{22--23}
  \answerto{Zwykle więcej niż jeden \dash{} $\val{p13.topo.count}$ dla
  sześciowierzchołkowego grafu w tym programie, spośród
  $\val{p13.topo.total}$ ustawień. Ważne jest to, że każdy z nich liczy te
  same wartości.}
\item Strona nie linkuje nigdzie. Co robi błądzenie losowe, gdy na niej
  wyląduje? \teachesat{31--32}
  \answerto{Zatrzymuje się, a masa prawdopodobieństwa opuszcza graf: po
  $\val{p13.drain.zeroat}$ krokach nie zostaje jej dokładnie nic. Po to jest
  współczynnik tłumienia i nie jest on sztuczką.}
\end{quiz}

% ============================================================
\section{Zbiór i relacja}
\label{sec:P13-what}
\index{graf!definicja}

\begin{fr}
Program~\ref{prog:F10} dał ci zbiór, a o grafach nie powiedział nic. Ten
program potrzebuje obok zbioru jednej rzeczy \dash{} zbioru par jego
elementów \dash{} a wszystko inne wynika z tych dwóch. W przeciwieństwie do
poprzedniego programu opiera się on na Części~III: do sekcji 3 cały temat
zamieni się w mnożenie macierzy.

Zacznij konkretnie. Sześć usług woła się nawzajem: brama woła
\code{auth} i \code{search}; \code{search} woła \code{rank} i
\code{store}; \code{rank} woła \code{store}; a \code{auth}, \code{rank} i
\code{store} piszą do \code{log}.

Zapomnij o rysunku, który właśnie kreślisz. Co najmniej trzeba zapisać, żeby
opisać to w całości?
\dotline
\nextframe
\end{fr}

\begin{fr}
\begin{ansblock}
Dwa zbiory: usługi oraz te ich pary, które są połączone.
\end{ansblock}

To cała definicja. \textbf{Graf} to zbiór $V$ \emph{wierzchołków} i zbiór $E$
ich par, czyli \emph{krawędzi}. Tutaj
$\lvert V \rvert = \val{p13.g.nv}$ oraz $\lvert E \rvert = \val{p13.g.ne}$.

Nic o położeniu, nic o przecinających się liniach, nic o wyglądzie. Rysunek
jest jednym przedstawieniem relacji, a takich są nieskończenie wiele; obiektem
jest relacja.

\begin{note}
To punkt Programu~\ref{prog:F10} o zbiorach, przybywający o poziom wyżej.
Zbiór odpowiada na pytanie \emph{czy to jest w środku}; graf odpowiada na
pytanie \emph{czy te dwa są połączone}. Wszystko w tym programie jest jednym z
tych dwóch pytań zadanym większej kolekcji.
\end{note}

Jednej rzeczy zdanie powyżej jednak nie powiedziało. \emph{Brama woła
\code{auth}} to nie to samo zdanie co \emph{\code{auth} woła bramę}. Co musi
się zmienić w definicji, żeby to wyrazić?
\nextframe
\end{fr}

\begin{fr}
\begin{ansblock}
Pary muszą być uporządkowane.
\end{ansblock}

Para \textbf{uporządkowana} $(u, v)$ czyni graf \emph{skierowanym}: krawędź
biegnie od $u$ do $v$ i nie z powrotem. Para nieuporządkowana $\{u, v\}$ czyni
go nieskierowanym, a wtedy relacja jest symetryczna z konstrukcji.

Ten jeden wybór to różnica między grafem wywołań, grafem importów i grafem
zależności z jednej strony \dash{} wszystkie skierowane \dash{} a grafem
znajomości czy grafem podobieństwa z drugiej. Warto decydować o tym wprost, bo
prawie każdy graf, jaki spotyka inżynier, jest skierowany, a prawie każdy
rysunek na tablicy zapomina o strzałkach.

\mermaidfig{p13-set-plus-relation}{Graf to dwa zbiory, a wybór par
uporządkowanych albo nie jest całym znaczeniem słowa
\enquote{skierowany}.}{dwa zbiory i relacja}
\end{fr}

\begin{fr}
Przeczytaj graf przez chwilę jako nieskierowany. \textbf{Stopień} wierzchołka
to liczba krawędzi przy nim, która sięga tu $\val{p13.g.degmax}$.

Dodaj wszystkie sześć stopni. Zrób to, zanim pójdziesz dalej, i powiedz, ile
ta suma musi wynosić w zależności od $\lvert E \rvert$.
\dotline
\nextframe
\end{fr}

\begin{fr}
\ans{$\val{p13.g.degsum}$, czyli $2\lvert E \rvert$}
A powód mieści się w jednym zdaniu: każda krawędź ma dwa końce, więc wnosi $1$
do stopnia przy każdym z nich, a dodanie stopni liczy każdą krawędź dokładnie
dwa razy.

\textbf{To podwójne zliczanie z~Programu~\ref{prog:P12} w nowym przebraniu.}
Nazywa się to zwykle lematem o uściskach dłoni, a skrypt sprawdza je na
$\val{p13.hand.trials}$ losowych grafach \dash{} co, jak
w~Programie~\ref{prog:P04}, jest testowaniem kodu, a nie szukaniem
kontrprzykładu, bo kontrprzykład obaliłby dowód.

\begin{note}
Bezpośredni wniosek warto mieć: skoro stopnie sumują się do liczby parzystej,
\textbf{wierzchołków o nieparzystym stopniu jest parzyście wiele.} Nie da się
mieć sieci, w której dokładnie jedna maszyna ma nieparzystą liczbę połączeń, a
taki argument o parzystości jest często najtańszym sprawdzeniem grafu, który
się właśnie zbudowało.
\end{note}
\end{fr}

\begin{fr}
Teraz drugi koniec skali. Graf, w którym połączona jest \emph{każda} para,
nazywa się pełnym. Jeśli ma $\val{p13.complete.n}$ wierzchołków, ile ma
krawędzi?
\nextframe
\end{fr}

\begin{fr}
\ans{$\val{p13.complete.e}$}
Czyli $\binom{n}{2}$, a to liczba par z~Programu~\ref{prog:F10}
i~Programu~\ref{prog:P12} co do cyfry \dash{} bo krawędź nieskierowana
\emph{jest} parą nieuporządkowaną, więc liczenie krawędzi w grafie pełnym to
liczenie par i nic więcej.

\begin{note}
Skrypt odczytuje zatwierdzoną liczbę par tamtego programu z jego pliku
wartości i przerywa budowanie, jeśli $\frac{n(n-1)}{2}$ jej nie odtworzy. Trzy
programy cytują teraz jedno obliczenie i żaden nie może się rozjechać.
\end{note}

Krawędzie dowolnego grafu o $n$ wierzchołkach są więc \emph{podzbiorem}
$\binom{n}{2}$ par, a faktycznie obecny ułamek to jego \textbf{gęstość}. To
słowo za chwilę rozstrzygnie, jak tę rzecz przechowywać.
\end{fr}

\begin{fr}
Jedno miejsce, w którym to nie jest analogia. Program~\ref{prog:P12}
deduplikował korpus, porównując każdą parę dokumentów. Zrób z tego graf:
dokumenty są wierzchołkami, a dwa są połączone, gdy ich podobieństwo
przekracza próg.

Wszystko, co tamten program mówił o parach, jest teraz zdaniem o krawędziach.
Liczba porównań to liczba par, czyli graf pełny; liczba \emph{trafień} to
liczba krawędzi, czyli o wiele mniej; a grupa dokumentów będących swoimi
duplikatami to zbiór wierzchołków połączonych ze sobą i z niczym innym.

\begin{aibox}
Ten ostatni obiekt ma nazwę i spotkałeś go bez niej. \textbf{Spójna składowa}
to zbiór wierzchołków osiągalnych z siebie nawzajem i z niczego spoza, a
znalezienie składowych jest tym, jak deduplikuje się korpus: zostaw po jednym
dokumencie z każdej.

To jest dokładnie tak, gdy \enquote{jest duplikatem} jest przechodnie, a
dopasowanie prawie-duplikatów zwykle nie jest. $A$ może pasować do $B$, a $B$
do $C$, podczas gdy $A$ i $C$ nie mają nic wspólnego, więc składowa łączy w
jedno dokumenty, z których żadna para nie pasowała \dash{} i tak przebieg
deduplikacji po cichu usuwa cały temat.

Nazwanie tego ma znaczenie, bo zmienia to, czego szukasz. \enquote{Które pary
pasują} to pytanie o $\binom{n}{2}$ rzeczy. \enquote{Jakie są składowe} to
pytanie o $n$ rzeczy, a na to jest standardowy algorytm, który biblioteka do
grafów już ma.
\end{aibox}
\end{fr}

% ============================================================
\section{Dwa zapisy, dwa koszty}
\label{sec:P13-encodings}
\index{graf!lista sąsiedztwa}

\begin{fr}
Relację da się zapisać na dwa sposoby i nie są one wymienne.

\textbf{Macierz sąsiedztwa} $A$ jest $n \times n$, z $A_{ij} = 1$, gdy istnieje
krawędź z $i$ do $j$, i $0$ w przeciwnym razie. \textbf{Lista sąsiedztwa}
trzyma dla każdego wierzchołka listę tych, na które wskazuje.

Przeczytaj ten graf jako nieskierowany, tak jak deduplikowany korpus powyżej:
krawędź łączy oba swoje końce, zamiast biec od jednego do drugiego, więc siedzi
na liście każdego z nich, a macierz wychodzi symetryczna. Definicje są ogólne,
a sekcje o DAG-ach i o błądzeniu losowym potrzebują ich właśnie takich; odczyt
zmienia nie to, czym jest kodowanie, ale ile miejsca zajmuje w nim krawędź.

Weź graf o $\val{p13.big.n}$ wierzchołkach i $\val{p13.big.m}$ krawędziach, co
jest zwykłym rozmiarem dla sieci cytowań albo grafu społecznego. Ile komórek
ma macierz?
\dotline
\nextframe
\end{fr}

\begin{fr}
\ans{$\val{p13.big.cells}$}
Milion do kwadratu. Przy jednym bajcie na komórkę \dash{} co już jest
absurdem, bo komórka mieści jeden bit informacji \dash{} daje to
$\val{p13.big.tb}$ terabajt na przechowanie $\val{p13.big.m}$ krawędzi.

Lista potrzebuje jednej głowy na wierzchołek i dwóch końców na krawędź:
$\val{p13.big.slots}$ miejsc, czyli $\val{p13.big.ratio}$ razy mniej.

\begin{trapbox}
\textbf{Macierz nie jest definicją, a lista nie jest implementacją.} Oba są
zapisami tej samej relacji i żaden nie jest bardziej prawdziwy.

Macierz wygląda na kanoniczną dlatego, że w niej pisze się matematykę, a lista
na praktyczną dlatego, że się mieści. Uczciwe zdanie brzmi: odpowiadają tanio
na różne pytania, a gęstość rozstrzyga, które pytanie będziesz zadawać. Tutaj
gęstość wynosi $\val{p13.big.density}$ \dash{} dwie krawędzie na każde sto
tysięcy możliwych \dash{} i \textbf{prawie każdy graf, jaki spotyka inżynier,
jest tak samo rzadki}, dlatego prawie każda biblioteka do grafów przyjmuje
listę jako domyślną.
\end{trapbox}
\end{fr}

\begin{fr}
Operacje różnią się równie ostro jak pamięć i warto ustalić, w którą stronę.

Dwa pytania. \emph{Czy istnieje krawędź z $u$ do $v$?} oraz \emph{jacy są
wszyscy sąsiedzi $u$?} Dla każdego powiedz, który zapis odpowiada szybciej i
dlaczego.
\dotline
\nextframe
\end{fr}

\begin{fr}
\begin{ansblock}
Macierz odpowiada na pierwsze w jednym kroku, a na drugie w $n$; lista
odpowiada na drugie w tylu krokach, ilu $u$ ma sąsiadów, a na pierwsze w
tyluż.
\end{ansblock}

Macierz ma odpowiedź na \emph{czy istnieje krawędź} pod znanym adresem, więc
jest to jedno odczytanie, jakkolwiek graf wygląda. Zapytanie jej o sąsiadów
oznacza przeczytanie całego wiersza, w większości zer.

Lista ma sąsiadów już zebranych, więc to pytanie kosztuje tyle, ile warta jest
odpowiedź. Zapytanie jej \emph{czy istnieje krawędź} oznacza przeszukanie
listy, co kosztuje tyle, co jej przeczytanie.

\begin{note}
Program~\ref{prog:P03} jest programem od mówienia tego w notacji $O$, a ten
celowo tego nie robi. \emph{To odczytanie czyta wiersz miliona komórek} jest
zdaniem użyteczniejszym na przeglądzie projektu niż \emph{to odczytanie jest
liniowe}, a dostępne bez żadnej notacji.
\end{note}
\end{fr}

\begin{fr}
Lista wygrywa więc na pamięci, wygrywa na sąsiadach i \emph{przegrywa} trzecie
pytanie: macierz odpowiada na \emph{czy istnieje krawędź z $u$ do $v$} w jednym
kroku, a lista przechodzi sąsiadów $u$ \dash{} w skali tej sekcji średnio
$\val{p13.big.avgdeg}$ z nich. To jedyna wygrana macierzy, a płacenie za nią
całymi wierszami jest powodem, dla którego nikt jej nie kupuje. Macierz byłaby
dziwną rzeczą do trzymania \dash{} poza jedną własnością, która nie ma nic
wspólnego z odczytywaniem czegokolwiek.

\textbf{Macierz można pomnożyć.} Następna sekcja jest o tym, co to daje, i to
jest powód, dla którego macierz sąsiedztwa przetrwała w temacie, w którym jest
$\val{p13.big.ratio}$ razy za duża, by ją przechować.
\end{fr}

% ============================================================
\section{Spacery to potęgi macierzy}
\label{sec:P13-walks}
\index{graf!spacer}

\begin{fr}
\textbf{Spacer} długości $k$ to ciąg $k$ krawędzi ułożonych koniec przy końcu.
Może odwiedzić wierzchołek ponownie; spacer, który tego nie robi, to
\emph{ścieżka}.

To rozróżnienie nie jest formalnością. Potęga macierzy liczy każdy ciąg
krawędzi, który układa się koniec przy końcu, a w iloczynie macierzy nie ma
nigdzie arytmetyki, która mogłaby pominąć ciągi wracające do już użytego
wierzchołka \dash{} iloczyn nie pamięta, gdzie był. Wszystko, co ta sekcja
dowodzi, dotyczy więc spacerów z samej konstrukcji, a liczenie
\emph{ścieżek} to problem istotnie trudniejszy, na który żadna potęga $A$
nie odpowiada.

A pytanie o spacery to prawie zawsze to pytanie, które masz. Pytanie, czy
zmiana tutaj może dosięgnąć usługi oddalonej o $k$ kroków, albo jaką część
grafu sygnał obejmie po $k$ rundach, jest pytaniem o ciągi krawędzi ułożone
koniec przy końcu i o nic więcej: trasa, która zawraca, zmarnowała krok, a i
tak dociera. Pytanie o ścieżki \dash{} ile jest różnych tras prostych
\dash{} zadaje się rzadko i odpowiada na nie drogo, więc tania maszyneria,
o której jest ta sekcja, celuje w to pytanie, które się pojawia.

Teraz pytanie, na którym stoi cała sekcja. Wypisz wyraz $(i,j)$ macierzy
$A^{2}$ jako sumę, używając wyłącznie definicji mnożenia macierzy:
\[ (A^{2})_{ij} = \sum_{k} A_{ik}A_{kj} \]
i powiedz, czym jest każdy wyraz tej sumy, skoro każdy wyraz $A$ to $0$ albo
$1$.
\dotline
\nextframe
\end{fr}

\begin{fr}
\begin{ansblock}
Każdy wyraz to $1$ dokładnie wtedy, gdy $i$ wskazuje na $k$ \emph{oraz} $k$
wskazuje na $j$, a $0$ w przeciwnym razie. Suma liczy więc wierzchołki
pośrednie \dash{} czyli spacery długości dwa z $i$ do $j$.
\end{ansblock}

\textbf{$(A^{k})_{ij}$ to liczba spacerów długości $k$ z $i$ do $j$}, a
indukcja od dwóch do $k$ to znowu to samo zdanie.

Zauważ, skąd to się wzięło. Program~\ref{prog:P06} dowodził, że iloczyn
macierzy \emph{jest} złożeniem i że wiersz-razy-kolumna to arytmetyka, która z
tego wypada. Spacer o dwóch krokach jest złożeniem dwóch kroków, więc liczba
spacerów jest iloczynem macierzy; twierdzenie to ten argument czytany w drugą
stronę i nie trzeba było do niego żadnej nowej maszynerii.

\begin{note}
Skrypt to sprawdza, zamiast wyprowadzać dwa razy: dla każdej pary
wierzchołków i każdej długości do $\val{p13.walk.maxk}$ wypisuje spacery
siłowo i wymaga, by ich liczba równała się wyrazowi macierzy \dash{}
$\val{p13.walk.checks}$ porównań, wszystkie dokładne na liczbach całkowitych.
Wypisanie nie może się mylić co do tego, czym jest spacer.
\end{note}
\end{fr}

\begin{fr}
Spróbuj na usługach. Ile spacerów długości $\val{p13.walk.k}$ biegnie od bramy
do \code{log}?
\nextframe
\end{fr}

\begin{fr}
\ans{$\val{p13.walk.count}$}
Dwie drogi przez \code{search}: przez \code{rank} i przez \code{store}.
Odczytanie ich z $A^{3}$ nie wymaga ani szukania, ani przechodzenia grafu, a
to samo mnożenie odpowiada na to pytanie dla wszystkich par naraz.

Zanim odpowiesz na następne, zauważ dwie rzeczy w tym, jak jest postawione.
Przekątna to ta część $A^{2}$, na którą nikt nie patrzy, bo \enquote{spacery
z $i$ z powrotem do $i$} brzmią jak trywialność \dash{} a twierdzenie
obejmuje ją bez wyjątku, więc cokolwiek ten wyraz mówi, jest wymuszone, a
nie wybrane. I słowo \emph{nieskierowanym} jest nośne: odpowiedź zmienia
się, gdy je pominąć, więc ustal, co to słowo kupuje, zanim się zdecydujesz.

Oto drugie czytanie tego samego twierdzenia, warte chwili, bo tu ludzie się
potykają. W grafie \emph{nieskierowanym} czym jest wyraz z przekątnej
$(A^{2})_{ii}$?
\dotline
\nextframe
\end{fr}

\begin{fr}
\begin{ansblock}
Stopniem $i$: spacery długości dwa z $i$ do siebie to dokładnie \enquote{krok
do sąsiada, krok z powrotem}.
\end{ansblock}

\begin{trapbox}
\textbf{Spacer to nie ścieżka, a przekątna jest miejscem, gdzie różnica
widać.} Wyjście i powrót to całkiem porządny spacer długości dwa i jest ich
tyle, ilu wierzchołek ma sąsiadów; nie jest to ścieżka i z pewnością nie jest
to cykl.

To pułapka w zdaniu \enquote{$A^{k}$ liczy trasy}, powiedzianym za szybko.
Jeśli chcesz ścieżek \dash{} tras, które nie powtarzają wierzchołka \dash{}
potęgi macierzy ich nie dają, a problem zliczania jest naprawdę trudniejszy.
Potęgi odpowiadają na to pytanie, na które odpowiadają, czyli na to o spacery,
i właśnie to pytanie zadaje przekazywanie komunikatów.
\end{trapbox}

\mermaidfig{p13-walks-are-powers}{Jeden krok to macierz, $k$ kroków to jej
$k$-ta potęga, a przekazywanie komunikatów to właśnie to mnożenie.}{kroki,
potęgi, komunikaty}
\end{fr}

\begin{fr}
To ostatnie zdanie jest zapłatą, więc oto ono w całości.

Grafowa sieć neuronowa daje każdemu wierzchołkowi wektor cech, układa je jako
wiersze macierzy $X$, a jedna warstwa liczy coś o kształcie
\[ X' = \sigma((A + I)XW) \]
Czytaj trzy czynniki od prawej: $W$ miesza własne cechy wierzchołka, $A + I$
zastępuje wiersz każdego wierzchołka sumą wierszy jego sąsiadów oraz jego
własnego, a $\sigma$ jest nieliniowością, bez której Program~\ref{prog:P06}
pokazał, że złożenie by się zapadło.

To $I$ jest pętlą własną, którą taka sieć dodaje przed agregacją, i nie jest
ozdobą: przy samym $A$ nowy wiersz wierzchołka powstaje z jego sąsiadów, a nie
z niego samego, więc warstwa zastępowałaby to, co wierzchołek wie, zamiast do
tego dokładać.

\textbf{Przekazywanie komunikatów jest mnożeniem przez macierz sąsiedztwa.}
Nie jest do niego podobne i nie jest tak zaimplementowane dla wygody; zdanie
\enquote{każdy wierzchołek zbiera od sąsiadów} i wyrażenie $AX$ to jedno i to
samo stwierdzenie.

Zatem: po $k$ takich warstwach, które wierzchołki mogły wpłynąć na wiersz
danego wierzchołka?
\dotline
\nextframe
\end{fr}

\begin{fr}
\ans{Dokładnie te w odległości $k$ od niego}
Bo $k$ warstw stosuje $A + I$ $k$ razy, a twierdzenie o liczeniu spacerów mówi,
czym są niezerowe wyrazy potęgi macierzy: to pary połączone spacerem długości
\textbf{dokładnie} $k$. Tak twierdzenie przybywa jako zdanie o architekturze
\dash{} a $I$ zamienia \emph{dokładnie} na \emph{w odległości}, bo krok zużyty
na pętlę własną stoi w miejscu, więc spacer dowolnej długości do $k$ zostaje
dopełniony do $k$ i policzony.

Zmierzone na usługach: przy jednym skoku brama sięga
$\val{p13.reach.one}$ wierzchołków, a przy $\val{p13.hops}$ skokach sięga
wszystkich $\val{p13.reach}$. Skrypt sprawdza ten zasięg wobec przeszukiwania
wszerz, które nie wie nic o macierzach, dla każdego wierzchołka i każdej
głębokości do $\val{p13.walk.maxk}$.

\begin{aibox}
\textbf{Głębokość w sieci przekazującej komunikaty to zasięg, a nie
pojemność}, i to prawdziwa różnica wobec każdej innej architektury w tej
książce. Dodanie warstwy do sieci gęstej daje jej więcej do liczenia; dodanie
warstwy do sieci grafowej daje każdemu wierzchołkowi szersze sąsiedztwo,
niezależnie od tego, czy było chciane.

Pytanie \enquote{ile warstw} ma więc odpowiedź płynącą z grafu, a nie z
budżetu: chodzi o to, jak daleko informacja naprawdę musi przebyć. W grafie o
średnicy $\val{p13.hops}$ czwarta warstwa nie dodaje do zasięgu nic nowego, a
w grafie typu małego świata dwie albo trzy warstwy dotykają już większości
wierzchołków.
\end{aibox}

\begin{warning}
Idący za tym tryb awarii \dash{} że przy dostatecznej liczbie warstw każdy
wierzchołek zbiera niemal z całego grafu, a reprezentacje zbiegają do siebie
\dash{} nazywa się zwykle nadmiernym wygładzaniem i \textbf{ta książka go nie
zmierzyła}. Ustalony jest tu zasięg, który jest faktem o $A^{k}$ i nie
potrzebuje eksperymentu. Czy i kiedy reprezentacje się zapadają, to twierdzenie
o wytrenowanych sieciach i takiego eksperymentu by wymagało.
\end{warning}
\end{fr}

% ============================================================
\section{DAG i porządek, który czyni obliczenie dobrze określonym}
\label{sec:P13-dag}
\index{graf!skierowany acykliczny}

\begin{fr}
Wróć do grafu wywołań i do własności, której jeszcze nie użyto: nie da się
wrócić tam, skąd się wyszło.

\textbf{Cykl} to spacer wracający do wierzchołka startowego bez powtarzania
krawędzi i bez przechodzenia dwukrotnie przez żaden inny wierzchołek. Graf
skierowany bez cykli jest \textbf{acykliczny}, a skierowany graf acykliczny to
\emph{DAG}.

Oba zastrzeżenia są potrzebne i wykluczają co innego. Każdy graf
nieskierowany mający choć jedną krawędź ma spacer, który wychodzi z
wierzchołka i wraca \dash{} owo \enquote{tam i z powrotem} z pułapki \dash{}
więc bez zastrzeżenia o krawędziach nic nie byłoby acykliczne, a słowo nie
znaczyłoby nic. W grafie skierowanym pracuje ono mniej, bo przejście krawędzią
z powrotem jest możliwe tylko wtedy, gdy taka krawędź istnieje; \emph{to} jest
powód, dla którego acykliczność orzeka się o grafach skierowanych, a o
nieskierowanych prawie nigdy. Zastrzeżenie o wierzchołkach wyklucza to, czego
nie wyklucza tamto: dwie pętle złączone w jednym punkcie, przebyte jako jedna
zamknięta trasa. Nie powtarza ona żadnej krawędzi, więc samo zastrzeżenie o
krawędziach by ją przepuściło, a jest to para cykli, a nie jeden.

Załóżmy, że musisz uruchomić te sześć usług w jakiejś kolejności, każdą po
wszystkim, od czego zależy. Co daje ci acykliczność?
\dotline
\nextframe
\end{fr}

\begin{fr}
\begin{ansblock}
To, że taka kolejność w ogóle istnieje.
\end{ansblock}

Kolejność, w której każdy wierzchołek stoi za wszystkim, co na niego wskazuje,
to \textbf{porządek topologiczny}, a twierdzenie mówi, że graf skierowany ma go
\emph{wtedy i tylko wtedy}, gdy jest acykliczny. Jedna strona jest łatwa i
warta zrobienia w głowie: gdyby był cykl, każdy jego wierzchołek musiałby stać
za następnym, a skończona lista tego nie potrafi.

Druga strona jest tym, co czyni słowo \emph{acykliczny} wartym powiedzenia, i
ta książka jej nie dowodzi \dash{} Program~\ref{prog:P14} jest tam, gdzie
czytania i formułowania twierdzeń uczy się porządnie.

Teraz je policz. Nasz graf ma $\val{p13.g.nv}$ wierzchołków, więc ustawień
jest $\val{p13.topo.total}$. Ile z nich to porządki topologiczne?
\nextframe
\end{fr}

\begin{fr}
\ans{$\val{p13.topo.count}$}
Cztery, spośród $\val{p13.topo.total}$ \dash{} ułamek
$\val{p13.topo.frac}$. Krawędzie są mocnym ograniczeniem i większość ustawień
łamie co najmniej jedno z nich.

Skrypt znajduje je, sprawdzając każdą permutację, a
Program~\ref{prog:P12} wycenia to na $n!$, co przybywa tu jako koszt: przy
sześciu jest dobrze, a
przy dwudziestu beznadziejnie, i prawdziwy planista używa algorytmu, a nie
przeszukania. Ważne jest tu nie to, jak się je znajduje, lecz to, że jest ich
więcej niż jeden.

Co prowadzi do pytania, dla którego ta sekcja istnieje. Te
$\val{p13.topo.count}$ porządki uruchamiają usługi w naprawdę różnych
sekwencjach. Czy liczą to samo?
\dotline
\nextframe
\end{fr}

\begin{fr}
\ans{Tak, zawsze}
Daj każdej usłudze wartość \dash{} powiedzmy o jeden większą od największej
wartości wśród usług, które ją wołają \dash{} a każdy porządek topologiczny
wyprodukuje identyczne wartości dla wszystkich sześciu. Nie zwykle, nie prawie:
identyczne, a skrypt to stwierdza dla wszystkich $\val{p13.topo.count}$.

\transcript{p13-order-and-answer}

Powód jest dość krótki, by go utrzymać. Wartość wierzchołka zależy tylko od
jego rodziców; porządek topologiczny liczy każdego rodzica przed
wierzchołkiem; więc jakikolwiek porządek wybrać, każdy wierzchołek widzi te
same wartości rodziców, gdy przychodzi jego kolej. Nic innego z porządku do
niego nie sięga.

\mermaidfig{p13-order-not-answer}{Wiele porządków, jedna odpowiedź \dash{} i to
pozwala planiście wybrać dowolny.}{porządek wolny, odpowiedź nie}
\end{fr}

\begin{fr}
\begin{aibox}
To cały powód, dla którego cztery na pozór niezwiązane rzeczy są jednym
obiektem.

\textbf{System budowania} ma pliki zależne od plików, uruchamia je w porządku
topologicznym i wolno mu zrównoleglić te części, które od siebie nie zależą
\dash{} czyli dokładnie tę swobodę, którą reprezentuje
$\val{p13.topo.count}$ porządków. \textbf{Przepływ agenta} ma kroki zależne od
kroków. \textbf{Przejście w przód sieci neuronowej} ma tensory zależne od
tensorów. \textbf{Plan zapytania} ma operatory zależne od operatorów.

Wszystkie cztery są DAG-ami, wszystkie cztery liczy się w porządku
topologicznym i wszystkie cztery są deterministyczne z tego samego powodu. Gdy
framework nazywa to, co buduje, \emph{grafem obliczeń}, nie jest to metafora.
\end{aibox}

Program~\ref{prog:P16} jest tam, gdzie tego grafu się używa, a nie definiuje:
przejście wstecz przechodzi ten sam DAG w odwrotnym porządku topologicznym, co
czyni \code{loss.backward()} dobrze określonym i dlatego aktywacje muszą tam
jeszcze być, gdy ono dotrze.
\end{fr}

\begin{fr}
A przypadek negatywny zasługuje na własną ramkę, bo to ten, który się spotyka.

Dodaj jedną krawędź do grafu wywołań \dash{} niech \code{log} woła bramę
\dash{} i zapytaj znowu o porządki topologiczne. Listing powyżej pokazuje już
odpowiedź: pustą listę.

\begin{trapbox}
\textbf{Cykl nie czyni obliczenia trudnym, tylko nieokreślonym.} Nie ma
kolejności, w której uruchomić usługi, bo każda musiałaby iść po następnej.
System budowania zgłasza cykl zależności i staje; nie jest nieuprzejmy, tam
naprawdę nie ma czego robić.

Dwa wyjścia warto rozpoznawać. \emph{Rozerwij cykl} \dash{} znajdź krawędź,
której nie powinno być, co w grafie wywołań jest zwykle złamaniem warstw. Albo
\emph{rozwiń go}: sieć rekurencyjna to cykl w architekturze i DAG w obliczeniu,
bo te same wagi stosuje się w każdym z ustalonej liczby kroków, a te kroki są
osobnymi wierzchołkami. Rozwinięcie jest tym, co daje trybowi wstecznemu
porządek topologiczny do pracy, i dlatego pamięć rośnie z liczbą
rozwiniętych kroków.
\end{trapbox}
\end{fr}

% ============================================================
\section{Gdzie osiada błądzenie losowe}
\label{sec:P13-walks-random}
\index{graf!błądzenie losowe}
\index{graf!PageRank}

\begin{fr}
Ostatni obiekt, i kładzie on maszynerię Części~III na grafie.

Zamiast pytać, które wierzchołki są osiągalne, idź za krawędziami losowo: z
wierzchołka wybierz jednostajnie jedną z krawędzi wychodzących i idź. Macierzą
tych prawdopodobieństw jest $P$, gdzie $P_{ij}$ to szansa kroku z $i$ do $j$
\dash{} macierz sąsiedztwa z każdym wierszem podzielonym przez jego własną
sumę, co wymaga, by ta suma była dodatnia. Wierzchołek bez krawędzi
wychodzących nie ma czego dzielić i jego wiersz zostaje taki, jaki był: same
zera.

Ile musi wynosić suma każdego wiersza $P$ i dlaczego?
\dotline
\nextframe
\end{fr}

\begin{fr}
\ans{$1$, bo błądzenie gdzieś idzie}
Każdy wiersz jest rozkładem prawdopodobieństwa nad tym, gdzie wyląduje
następny krok, a podzielenie wiersza jedynek przez liczbę jedynek jest
dokładnie tym, co czyni go jednością.

Teraz połóż rozkład $p$ na wierzchołkach \dash{} szansę bycia w każdym \dash{}
i zrób krok. Nowym rozkładem jest $P\T p$: szansa dotarcia do $j$ to suma po
$i$ z \emph{bycia w $i$} razy \emph{kroku z $i$ do $j$}, czyli iloczyn
macierzy przez wektor i nic więcej.

Zacznij jednostajnie i iteruj. Powiedz, czego się spodziewasz po $p$.
\nextframe
\end{fr}

\begin{fr}
\begin{ansblock}
Osiada: po dostatecznej liczbie kroków rozkład przestaje się zmieniać.
\end{ansblock}

Zmierzone na czterostronicowym grafie linków, po $\val{p13.pi.steps}$ krokach:
$\val{p13.pi.0}$, $\val{p13.pi.1}$, $\val{p13.pi.2}$, $\val{p13.pi.3}$.
Wartości dokładne to
$\frac{\val{p13.exact.num.0}}{\val{p13.exact.den.0}}$,
$\frac{\val{p13.exact.num.1}}{\val{p13.exact.den.1}}$,
$\frac{\val{p13.exact.num.2}}{\val{p13.exact.den.2}}$ oraz
$\frac{\val{p13.exact.num.3}}{\val{p13.exact.den.3}}$, rozwiązane na liczbach
wymiernych, a nie iterowane, a iteracja zgadza się z nimi lepiej niż
$10^{-\val{p13.pi.bound}}$.

Rozkład, który krok zostawia w spokoju, nazywa się \textbf{stacjonarnym}.
Zapisz ten warunek jako równanie w $P\T$ i $p$ i powiedz, na co patrzysz.
\dotline
\nextframe
\end{fr}

\begin{fr}
\ans{$P\T p = p$, czyli wektor własny dla wartości własnej $1$}
Program~\ref{prog:P10} zdefiniował dokładnie to: kierunek, którego macierz nie
obraca, tylko rozciąga, a tutaj współczynnik rozciągnięcia wynosi $1$.

Rozkład stacjonarny błądzenia losowego jest więc rachunkiem na wektorze
własnym, a odczyt, który daje tamten program, stosuje się bez zmian
\dash{} wielokrotne stosowanie macierzy mnoży każdy kierunek własny przez
jego własny współczynnik podniesiony do tej potęgi, więc ten, którego
współczynnik zostawia wektor bez zmian, jest tym, na czym osiada długie
błądzenie. Ani iteracyjnej metody jego szukania, ani tempa, w jakim taka
metoda by osiadała, nie ma w tamtym programie ani w tej książce.

\begin{note}
Skrypt nie iteruje z nadzieją. Rozwiązuje $(P\T - I)p = 0$ przy
$\sum p = 1$ dokładnie na liczbach wymiernych, tą samą eliminacją Gaussa,
którą napisał Program~\ref{prog:P04}, a Programy~\ref{prog:P08}
i~\ref{prog:P09} wykorzystały, a potem stwierdza $P\T p = p$ bez żadnej
tolerancji. \enquote{Stacjonarny} to twierdzenie o równości, a porównanie
zmiennoprzecinkowe uczyniłoby z niego twierdzenie o progu.
\end{note}
\end{fr}

\begin{fr}
To jest \textbf{PageRank} i definicja jest całą jego treścią: rangą strony
jest jej udział w czasie, który spędza tam nieskończone błądzenie losowe.
Strona jest ważna, bo linkują do niej ważne strony, co brzmi na okrężne, a
jest równaniem na wektor własny.

To rozróżnienie warto mieć w słowach, bo \enquote{okrężne} to pierwszy
zarzut, jaki ktokolwiek podnosi. Okrężna \emph{definicja} tłumaczy pojęcie
przez nie samo i nie tłumaczy niczego. Zapis $P\T p = p$ nie jest definicją,
tylko \emph{warunkiem}: nie mówi nic o tym, czym ranking jest, a zamiast tego
żąda, by \dash{} czymkolwiek jest \dash{} zgadzał się ze sobą po jednym kroku
błądzenia. Czy jakikolwiek taki ranking w ogóle istnieje, jest wtedy
prawdziwym pytaniem o prawdziwej odpowiedzi, a odpowiedzią jest wektor własny,
który istnieje, a przy warunku z kolejnych dwóch ramek jest jedyny.

Ale błądzenie tak opisane psuje się na grafie, jaki sieć z pewnością zawiera.
Co się dzieje, gdy wyląduje na stronie bez żadnych linków wychodzących?
\dotline
\nextframe
\end{fr}

\begin{fr}
\begin{ansblock}
Nie ma dokąd pójść, a masa prawdopodobieństwa opuszcza graf.
\end{ansblock}

Ten wiersz $P$ jest samymi zerami, więc nie jest rozkładem, a iteracja nie
zachowuje sumy. Zmierzone na trzystronicowym grafie z jedną ślepą uliczką:
masa pozostała po kolejnych krokach to
$\frac{\val{p13.drain.num1}}{\val{p13.drain.den1}}$, potem
$\frac{\val{p13.drain.num2}}{\val{p13.drain.den2}}$, a potem
\textbf{dokładnie zero} do kroku $\val{p13.drain.zeroat}$.

Nie prawie zero. Zero, na liczbach wymiernych, ze wszystkim, co odeszło.
\end{fr}

\begin{fr}
Naprawa mieści się w jednej linijce i zwykle opisuje się ją jako sztuczkę,
którą nie jest.

Z prawdopodobieństwem $\val{p13.damp}$ idź za linkiem jak dotąd; z pozostałym
skocz do wierzchołka wybranego jednostajnie losowo. Strona bez linków skacze za
każdym razem. Wynikają z tego od razu dwie rzeczy i obie są sednem.

\textbf{Każdy wyraz macierzy jest teraz dodatni}, więc każda strona jest
osiągalna z każdej innej w jednym kroku, a błądzenie nie da się uwięzić w
podgrafie linkującym tylko do siebie. Oraz \textbf{nic nie wycieka}, bo każdy
wiersz znowu sumuje się do jedności. Zmierzone na tym samym trzystronicowym
grafie, tłumione błądzenie osiada na $\val{p13.damped.0}$,
$\val{p13.damped.1}$ i $\val{p13.damped.2}$, a ślepa uliczka trzyma największy
udział, zamiast pochłaniać wszystko.

\begin{note}
Wartość $\val{p13.damp}$ jest umowna, a nie wyprowadzona, i ta książka nie
będzie udawać inaczej. \emph{Wyprowadzone} jest to, że tłumienie musi tam być:
bez niego punkt stały albo nie istnieje, albo nie jest jedyny, a z nim macierz
jest dodatnia, o co prosi twierdzenie Perrona--Frobeniusa.
\end{note}
\end{fr}

\begin{fr}
Jedna uwaga na koniec, i o to tak naprawdę w tym programie chodziło.

Zadano jednemu obiektowi trzy pytania i każde zamieniło się w arytmetykę już
obecną w książce. \emph{Co może sięgnąć czego} stało się potęgą macierzy. \emph{W
jakim porządku wolno to liczyć} stało się własnością, której sprawdzenie nic
nie kosztuje, a kupuje determinizm. \emph{Gdzie osiada błądzenie losowe} stało
się wektorem własnym.

\textbf{Żadne z tych trzech nie potrzebowało rysunku}, a rysunek jest tym, po
co większość ludzi sięga najpierw. Graf narysowany na tablicy jest
przedstawieniem relacji; to relacja jest tym, o czym mówi arytmetyka, i tylko
ona przetrwa bycie szeroką na milion wierzchołków.
\end{fr}

% ============================================================
\begin{summarybox}
\sumitem{1--3}{\result{\textbf{Graf to zbiór wierzchołków i zbiór ich par}, i nic
  więcej \dash{} bez położenia, bez rysunku. \textbf{Pary uporządkowane czynią
  go skierowanym}, co mają graf wywołań, graf importów i graf zależności, a
  szkic na tablicy zwykle gubi}.}
\sumitem{4--5}{\result{\textbf{Stopnie sumują się do podwojonej liczby krawędzi}, bo
  każda krawędź jest liczona przy każdym ze swoich dwóch końców:
  $\val{p13.g.degsum}$ przy $\val{p13.g.ne}$ krawędziach. Wierzchołków o
  nieparzystym stopniu jest więc parzyście wiele}.}
\sumitem{6--8}{\result{Graf pełny o $\val{p13.complete.n}$ wierzchołkach ma
  $\val{p13.complete.e}$ krawędzi, czyli liczbę par z~Programów~\ref{prog:F10}
  i~\ref{prog:P12} \dash{} krawędź nieskierowana \emph{jest} parą
  nieuporządkowaną. Deduplikacja to znajdowanie spójnych składowych}.}
\sumitem{9--10}{\result{\textbf{Macierz czy lista to pytanie o gęstość i o to, którą
  operację powtarzasz.} Przy $\val{p13.big.n}$ wierzchołkach i
  $\val{p13.big.m}$ krawędziach macierz ma $\val{p13.big.cells}$ komórek wobec
  $\val{p13.big.slots}$ miejsc, a gęstość wynosi $\val{p13.big.density}$}.}
\sumitem{11--13}{\result{Macierz odpowiada na \emph{czy istnieje krawędź} w jednym
  kroku, a lista na \emph{jacy są sąsiedzi} w tylu krokach, ilu jest sąsiadów.
  \textbf{Macierz zarabia na swoje tym, że da się ją pomnożyć}, a nie tym, że
  da się w niej wyszukać}.}
\sumitem{14--15}{\result{\textbf{$(A^{k})_{ij}$ liczy spacery długości $k$ z $i$ do
  $j$}, bo każdy wyraz $\sum_k A_{ik}A_{kj}$ to $1$ dokładnie wtedy, gdy oba
  kroki istnieją \dash{} złożenie z~Programu~\ref{prog:P06} robiące
  arytmetykę}.}
\sumitem{16--18}{\result{\textbf{Spacer to nie ścieżka}: $(A^{2})_{ii}$ to stopień
  $i$, liczący \enquote{tam i z powrotem}. Zliczanie ścieżek jest problemem
  trudniejszym i potęgi macierzy tego nie robią}.}
\sumitem{19--20}{\result{\textbf{Przekazywanie komunikatów jest mnożeniem przez
  macierz sąsiedztwa}, więc $k$ warstw sięga w promieniu $k$ skoków, gdy
  warstwa dodaje pętlę własną:
  $\val{p13.reach.one}$ wierzchołków przy jednym skoku i wszystkie
  $\val{p13.reach}$ przy $\val{p13.hops}$. \textbf{Głębokość to zasięg, a nie
  pojemność.}}}
\sumitem{21--23}{\result{\textbf{Graf skierowany ma porządek topologiczny wtedy i
  tylko wtedy, gdy jest acykliczny}, a zwykle ma ich więcej niż jeden:
  $\val{p13.topo.count}$ tutaj spośród $\val{p13.topo.total}$}.}
\sumitem{24--25}{\result{\textbf{Każdy porządek topologiczny liczy te same wartości},
  bo każdy wierzchołek zależy tylko od rodziców, których wszystkie porządki
  liczą wcześniej. Dlatego system budowania, przepływ agenta, plan zapytania i
  przejście w przód są jednym obiektem i dlatego planista może swobodnie
  zrównoleglać}.}
\sumitem{26}{\result{\textbf{Cykl czyni obliczenie nieokreślonym, a nie trudnym.}
  Rozerwij go albo rozwiń \dash{} czym jest sieć rekurencyjna i dlaczego jej
  pamięć rośnie z liczbą rozwiniętych kroków}.}
\sumitem{27--30}{\result{Krokiem błądzenia losowego jest $P\T p$, a
  \textbf{rozkład stacjonarny to wektor własny dla wartości własnej $1$}
  \dash{} Program~\ref{prog:P10} pod inną nazwą. To jest PageRank i sprawdza
  się go tutaj jako dokładną równość, a nie tolerancję}.}
\sumitem{31--34}{\result{\textbf{Wierzchołek bez krawędzi wychodzącej wysusza
  błądzenie}: po $\val{p13.drain.zeroat}$ krokach nie zostaje dokładnie nic
  masy. Tłumienie na $\val{p13.damp}$ czyni każdy wyraz dodatnim, a każdy
  wiersz sumującym się do jedności, i to kupuje jedyny punkt stały \dash{}
  stała jest umowna, potrzeba jej nie}.}
\end{summarybox}

\canyou

\begin{testexercises}
\item Graf skierowany ma $12$ wierzchołków i $30$ krawędzi. Ile wynosi suma
  stopni wyjściowych? A gdyby był nieskierowany, suma stopni?
  \answerto{$30$ tak czy inaczej dla stopni wyjściowych \dash{} każda krawędź
  skierowana wychodzi z dokładnie jednego wierzchołka. Nieskierowany daje
  $60$, dwa razy tyle co krawędzi, bo każda krawędź jest liczona przy obu
  końcach. Te dwie odpowiedzi różnią się o czynnik dwa, a mylenie ich to
  najczęstsze potknięcie w tym temacie.}
\item Masz $50\,000$ wierzchołków i $200\,000$ krawędzi. Podaj rozmiary
  macierzy i listy i powiedz, co byś przechował.
  \answerto{$50\,000^{2} = \num{2.5e9}$ komórek wobec
  $50\,000 + 400\,000 = \num{450000}$ miejsc, czyli około $5500$ razy mniej.
  Przechowaj listę; gęstość wynosi $\num{1.6e-4}$. Rzadką macierz trzymaj
  dodatkowo tylko wtedy, gdy zamierzasz przez nią mnożyć.}
\item $A$ jest macierzą sąsiedztwa grafu nieskierowanego. Ile wynosi suma
  wszystkich wyrazów $A$? A suma przekątnej $A^{2}$?
  \answerto{Obie to podwojona liczba krawędzi. Pierwsza, bo każda krawędź
  ustawia dwa wyrazy; druga, bo $(A^{2})_{ii}$ to stopień $i$, a stopnie
  sumują się do $2\lvert E \rvert$. Ta sama wielkość dwiema drogami, czyli
  znowu lemat o uściskach dłoni.}
\item Budowanie ma $8$ celów i jeden z nich, przechodnio, zależy od siebie. Co
  zrobi system budowania i jakie masz dwie możliwości?
  \answerto{Zgłosi cykl zależności i odmówi uruchomienia, bo nie ma
  kolejności, w której każdy cel stoi za tym, czego potrzebuje. Albo rozerwij
  cykl \dash{} zwykle złamanie warstw \dash{} albo rozwiń go w osobne kopie,
  co robi sieć rekurencyjna i dlaczego jej pamięć rośnie z liczbą kroków.}
\item Grafowa sieć neuronowa ma $6$ warstw, a jej graf ma średnicę $3$. Co
  szósta warstwa dodaje do zasięgu któregokolwiek wierzchołka?
  \answerto{Nic. Po $3$ warstwach każdy wierzchołek zbiera już z całej spójnej
  składowej, więc warstwy od czwartej do szóstej dodają parametry i obliczenia,
  a żadnej nowej informacji. Głębokość w sieci przekazującej komunikaty to
  zasięg, a zasięg kończy się na średnicy.}
\end{testexercises}

\begin{furtherproblems}
\item Pokaż, że wierzchołków o nieparzystym stopniu jest parzyście wiele.
  \answerto{Stopnie sumują się do $2\lvert E \rvert$, czyli do liczby
  parzystej. Wierzchołki o parzystym stopniu wnoszą sumę parzystą, więc te o
  nieparzystym też muszą \dash{} a suma liczb nieparzystych jest parzysta
  tylko wtedy, gdy jest ich parzyście wiele.}
\item $A$ ma wiersz samych zer. Co to mówi o grafie i co robi z $A^{k}$?
  \answerto{Ten wierzchołek nie ma krawędzi wychodzącej. Wiersz $i$ macierzy
  $A^{k}$ jest wtedy zerowy dla każdego $k$, więc żaden spacer żadnej długości
  z niego nie wychodzi \dash{} czyli ta sama ślepa uliczka, która wysusza
  błądzenie losowe w sekcji 5, widziana w macierzy sąsiedztwa zamiast w
  macierzy przejść.}
\item Dlaczego zliczanie ścieżek nie sprowadza się do potęgi macierzy tak, jak
  zliczanie spacerów?
  \answerto{Bo \enquote{nie powtarza wierzchołka} jest warunkiem na cały ciąg,
  a mnożenie macierzy łączy kroki, nie wiedząc, co było przed nimi. Suma
  $\sum_k A_{ik}A_{kj}$ nie potrafi wykluczyć wyrazów, w których $k$ równa się
  $i$ albo $j$, jeśli jej się tego nie powie, a przy długości trzy warunek
  dotyczy całej trasy, a nie jednego wyrazu.}
\item DAG ma $n$ wierzchołków i ani jednej krawędzi. Ile ma porządków
  topologicznych? A przy jednym łańcuchu $n-1$ krawędzi?
  \answerto{$n!$ bez krawędzi \dash{} kwalifikuje się każde ustawienie \dash{}
  i dokładnie $1$ dla łańcucha, który ustala porządek w pełni. Każdy DAG leży
  między tymi dwoma, a ta liczba mierzy, ile swobody ma planista.}
\item Współczynnik tłumienia opisuje się często jako \enquote{szansę, że
  surfujący się znudzi}. Co on faktycznie kupuje, matematycznie?
  \answerto{Dwie rzeczy. Każdy wyraz macierzy staje się dodatni, więc żaden
  zbiór stron nie uwięzi błądzenia, i to czyni punkt stały jedynym; oraz każdy
  wiersz sumuje się do jedności nawet dla strony bez linków, więc nic nie
  wycieka. Opowieść o nudzie to motywacja; tamte dwie własności są tym, czego
  potrzebuje twierdzenie.}
\end{furtherproblems}
